\section{Standardness outside characteristic two}
\label{sec:mesh-deformations}

We now complete the proof of Theorems~A and~B when $\charac k\ne2$. By
Corollary~\ref{cor:graded-standard-form}, it suffices to prove that
$A\cong A_0$, or equivalently that $\ind(\modu A)\cong\C$. In the graded
category $\C=\gr_\JJ\ind(\modu A)$, the chosen arrow representatives satisfy
the quadratic mesh relations \eqref{eq:scalar-mesh} exactly, whereas any
lifts of those arrows to $\ind(\modu A)$ satisfy the mesh relations modulo
$\JJ^3$. If the linear derivative map sending degree-$m$ perturbations of the
arrows to their degree-$(m+1)$ mesh errors is surjective for every $m\geq2$,
we can inductively adjust the chosen irreducible morphisms degree by degree
until the mesh relations hold on the nose in $\ind(\modu A)$.

In Subsection~\ref{subsec:mesh-derivative} we show that, in any
characteristic, the cokernel of this linearized mesh derivative localizes to
the $\tau$-periodic components of $\Gamma_A$. In
Appendix~\ref{app:periodic-mesh} we prove that the localized cokernel
vanishes whenever $\charac k\ne2$, provided the mesh category satisfies a
natural necessary condition for finite representation type that we call the
\emph{brick hypothesis}: every nonempty object quotient contains an object
with endomorphism ring $k$. Subsection~\ref{subsec:brick-hypothesis} verifies
that $\C$ satisfies this hypothesis (which rules out the small-period
components of types $E_6$, $E_7$, and $E_8$ that would otherwise have nonzero
cokernels in characteristics $3$ and $5$), and
Subsection~\ref{subsec:successive-correction} carries out the successive
arrow corrections.

\subsection{The derivative and periodic vertices}
\label{subsec:mesh-derivative}

Throughout this subsection, let $Q$ be an arbitrary finite translation quiver
with at most one arrow between each ordered pair of vertices, and let
$C=C(Q,c)$ be the scalar mesh category of
Definition~\ref{def:scalar-mesh-category} for a fixed system of nonzero
coefficients $c=(c_{x,a})$ (we reserve the calligraphic letter $\C$ for the
graded module category $\gr_\JJ\ind(\modu A)$).

For each integer $m\geq2$, define the homogeneous spaces of degree-$m$ arrow
perturbations and degree-$(m+1)$ mesh errors by
\[
 U_m(C)=\bigoplus_{a:u\to v}C_m(u,v),\qquad
 V_{m+1}(C)=\bigoplus_{x\text{ nonprojective}}C_{m+1}(\tau x,x),
\]
where $C_m$ denotes the homogeneous component of path length $m$. Substituting
$a\mapsto a+f_a$ with $f_a\in C_m(u,v)$ into the quadratic mesh relations
$r_x=\sum_{a:y\to x}c_{x,a}(\sigma a)a$ and taking the linear term in $f$
yields the \emph{mesh derivative}
\begin{equation}
 \begin{split}
 D_m:U_m(C)&\longrightarrow V_{m+1}(C),\\
 (D_mf)_x&=\sum_{a:y\to x}c_{x,a}
             \bigl(f_{\sigma a}a+(\sigma a)f_a\bigr).
 \end{split}
 \label{eq:mesh-derivative}
\end{equation}
Thus $\operatorname{coker}D_m(C)$ measures the degree-$(m+1)$ mesh errors
that cannot be canceled by a degree-$m$ modification of the arrows.

Let $T\subseteq Q_0$ be the set of $\tau$-periodic vertices of $Q$, so that
the full subquiver $Q_T$ is a stable translation quiver, and let
$C_T=C(Q_T,c|_T)$ be its scalar mesh category, which by
\eqref{eq:mesh-object-quotient} is the object quotient of $C$ obtained by
killing the identity morphisms outside $T$ (compare Riedtmann's stable-block
construction \cite{Riedtmann1984}). The following localization proposition is
used both in Appendix~\ref{app:periodic-mesh} (for $\charac k\ne2$) and in
Section~\ref{sec:positive-deformations} (for $\charac k=2$).

\begin{proposition}
\label{prop:periodic-localization}
Let $C=C(Q,c)$ be a scalar mesh category on a finite translation quiver $Q$
with at most one arrow between each ordered pair of vertices, and let $C_T$
be its object quotient on the set $T$ of $\tau$-periodic vertices. For every
$m\geq2$, the quotient functor $C\twoheadrightarrow C_T$ induces an
isomorphism
\[
 \operatorname{coker}D_m(C)\xrightarrow{\sim}
 \operatorname{coker}D_m(C_T).
\]
\end{proposition}
\begin{proof}
Before imposing the mesh relations of $C$, let $\widetilde U_m$ and
$\widetilde V_{m+1}$ be the free vector spaces spanned by literal paths in
$Q$, with basis $(b,p)$ for $\widetilde U_m$ (where $b:u\to v$ is an arrow
and $p:u\to v$ is a path of length $m$) and basis $(x,w)$ for
$\widetilde V_{m+1}$ (where $x$ is a nonprojective vertex and
$w:\tau x\to x$ is a path of length $m+1$). Formula~\eqref{eq:mesh-derivative}
defines a lifted linear map $\widetilde D_m:\widetilde U_m\to\widetilde V_{m+1}$.

In the matrix of $\widetilde D_m$ with respect to these path bases, each
column $(b,p)$ has at most two nonzero entries: if $v$ is nonprojective, the
occurrence of $b$ as the second arrow in $r_v$ contributes
$c_{v,b}(\sigma b)p$ to the row $(v,(\sigma b)p)$; if $u$ is noninjective,
writing $b=\sigma b'$ for $b':v\to\tau^{-1}u$, the occurrence of $b$ as the
first arrow in $r_{\tau^{-1}u}$ contributes $c_{\tau^{-1}u,b'}pb'$ to the row
$(\tau^{-1}u,pb')$. Conversely, each row $(x,w)$ with $w:\tau x\to x$ of
length $m+1$ is incident to exactly two columns (counted with multiplicity),
corresponding to the first and last arrows of $w$: writing $w=bp$ with first
arrow $b=\sigma a:\tau x\to y$ gives the incidence with column $(a,p)$, while
writing $w=pa$ with last arrow $a:y\to x$ gives the incidence with column
$(\sigma a,p)$, which links $(x,pa)$ to $(y,(\sigma^2 a)p)$ whenever $y$ is
nonprojective and has no second entry when $y$ is projective.

Viewing the rows $(x,w)$ as vertices and the two-entry columns as edges (with
one-entry columns as half-edges at exposed endpoints), the connected
components of this incidence graph are finite chains and finite cycles. Every
chain component ends at a one-entry column, so Gaussian elimination from that
exposed end places the entire chain component in $\im\widetilde D_m$. Along a
two-entry column, moving from $(x,pa)$ to $(y,(\sigma^2 a)p)$ rotates the
last arrow $a$ of $w$ to the front and applies $\tau=\sigma^2$; iterating this
rotation $d=m+1$ times gives
\[
 \operatorname{rot}^d(x,w)=(\tau x,\tau w).
\]
Consequently, a row $(x,w)$ lies on a cycle component if and only if every
vertex visited by $w$ admits arbitrarily many positive and negative powers of
$\tau$, which in the finite quiver $Q$ means precisely that the path $w$ lies
entirely inside $Q_T$. Moreover, because a vertex is in $T$ if and only if
its $\tau$-translate is in $T$, no column of $\widetilde D_m$ links a path
contained in $Q_T$ to a path passing outside $T$. The free row space
$\widetilde V_{m+1}$ therefore decomposes as a direct sum of the subspace of
paths contained in $Q_T$ and the subspace of paths passing outside $T$; the
latter is a sum of chain components and lies in $\im\widetilde D_m$, so
projection onto paths in $Q_T$ induces an isomorphism
$\operatorname{coker}\widetilde D_m(Q)\xrightarrow{\sim}\operatorname{coker}\widetilde D_m(Q_T)$.

Finally, let $R_V\subseteq\widetilde V_{m+1}$ be the subspace spanned by the
contextual mesh relations in $\widetilde V_{m+1}$, so that
\[
 \operatorname{coker}D_m(C)=\widetilde V_{m+1}/\bigl(R_V+\im\widetilde D_m\bigr).
\]
Because the projection onto paths contained in $Q_T$ kills every contextual
mesh relation passing outside $T$ and sends the remaining terms to the
contextual mesh relations of $C_T$, it induces the stated isomorphism
$\operatorname{coker}D_m(C)\xrightarrow{\sim}\operatorname{coker}D_m(C_T)$.
\end{proof}

\subsection{The restriction imposed by finite representation type}
\label{subsec:brick-hypothesis}

We say that a $k$-linear category $C$ satisfies the \emph{brick hypothesis}
if every nonempty object quotient of $C$ contains a \emph{brick}, that is, an
object whose endomorphism ring is $k$. Here, for any subset $\Sigma$ of
objects of $C$, the object quotient $C/[\Sigma]$ denotes the quotient of $C$
by the ideal of morphisms factoring through finite direct sums of objects in
$\Sigma$, with the objects in $\Sigma$ removed. Although the periodic
quotient $\C_T$ of $\C=\gr_\JJ\ind(\modu A)$ need not itself be the category of
indecomposable modules of an algebra, both $\C$ and $\C_T$ inherit the brick
hypothesis from $\modu A_0$.

\begin{lemma}
\label{lem:deletion-brick}
Let $\mathcal D$ be a full additive subcategory of a finite-length
$k$-linear abelian category, closed under images and direct summands, with
finite-dimensional Hom spaces and Krull--Schmidt decomposition. For any subset
$\Sigma\subseteq\operatorname{ob}(\ind\mathcal D)$, every nonempty object
quotient $(\ind\mathcal D)/[\Sigma]$ contains an object whose endomorphism
ring is a division algebra (and hence is equal to $k$ when $k$ is
algebraically closed).
\end{lemma}
\begin{proof}
Choose an object $X\in\operatorname{ob}(\ind\mathcal D)\setminus\Sigma$ of
minimal composition length. Every radical endomorphism $f\in\rad\End(X)$ has
a proper image $\im f$ of strictly smaller composition length. By hypothesis,
$\im f$ belongs to $\mathcal D$, and each of its indecomposable summands has
length smaller than that of $X$ and therefore belongs to $\Sigma$. Thus every
radical endomorphism of $X$ vanishes in $(\ind\mathcal D)/[\Sigma]$. On the
other hand, $1_X$ cannot factor through a finite direct sum of objects in
$\Sigma$ by the Krull--Schmidt property, so the endomorphism ring of $X$ in
$(\ind\mathcal D)/[\Sigma]$ is the residue division ring of the local ring
$\End(X)$.
\end{proof}

By Proposition~\ref{prop:projective-evaluation}(1), $\add\C\simeq\modu A_0$
is a finite-length abelian category, so Lemma~\ref{lem:deletion-brick} shows
that $\C$ satisfies the brick hypothesis. Since every object quotient of
$\C_T$ is also an object quotient of $\C$, the periodic mesh category $\C_T$
satisfies the brick hypothesis as well. Using Proposition~\ref{prop:periodic-tree-model}
and Proposition~\ref{prop:periodic-localization}, we prove the following
vanishing theorem in Appendix~\ref{app:mesh-calculations}.

\begin{proposition}
\label{prop:mesh-correction-spaces}
Let $C=C(Q,c)$ be a finite-dimensional scalar mesh category on a finite
translation quiver $Q$ with at most one arrow between each ordered pair of
vertices. If $\charac k\ne2$ and $C$ satisfies the brick hypothesis, then the
mesh derivative $D_m(C)$ is surjective for every $m\geq2$.
\end{proposition}

\subsection{Successive correction of irreducible representatives}
\label{subsec:successive-correction}

When $D_m$ is surjective for all $m\geq2$, a standard degree-by-degree
perturbation argument (compare Riedtmann \cite[Section~3.3]{Riedtmann1984})
shows that any filtered category whose associated graded category is $C$ is
already isomorphic to $C$.

\begin{proposition}
\label{prop:filtered-normal-form}
Let $C=kQ/(r_1,\ldots,r_t)$ be a finite-dimensional path-length graded
algebra over an arbitrary field $k$, where $Q$ is a finite quiver and each
relation $r_i$ is homogeneous of degree $2$. For each $q\geq3$, let
$D_{q-1}$ denote the linear map sending a tuple $(f_a)_a$ of homogeneous
elements $f_a\in C_{q-1}$ parallel to the arrows $a$ of $Q$ to the tuple in
$C_q$ obtained by substituting $a\mapsto a+f_a$ into $(r_1,\ldots,r_t)$ and
retaining the linear terms in $f$. If $D_{q-1}$ is surjective for every
$q\geq3$, then any finite-dimensional filtered algebra or category $E$
equipped with a graded isomorphism $\gr E\cong C$ is isomorphic to $C$ via a
filtered isomorphism inducing the given identification on $\gr E$.
\end{proposition}
\begin{proof}
Let $F^\bullet E$ be the finite decreasing filtration of $E$. Lift the vertex
idempotents of $C$ to a complete set of orthogonal idempotents of $E$ (taking
the identity morphisms when $E$ is a category), and choose initial lifts in
$F^1 E$ of the arrows $a$ in the corresponding corners of $E$. Their
quadratic relation errors $r_i$ then lie in $F^3 E$.

Suppose by induction that we have chosen arrow lifts whose relation errors
lie in $F^q E$ for some $q\geq3$. Passing to $F^q E/F^{q+1}E\cong C_q$, the
leading classes of these relation errors lie in the image of $D_{q-1}$ by
hypothesis. Subtracting a parallel correction of degree $q-1$ from each arrow
lift cancels the degree-$q$ part of each relation error, while the terms
quadratic in the correction have degree at least $2(q-1)>q$; thus the new
relation errors lie in $F^{q+1}E$. Because the filtration $F^\bullet E$ is
finite, this process terminates after finitely many steps, yielding arrow
lifts in $E$ that satisfy the relations $r_1,\ldots,r_t$ on the nose. The
induced homomorphism $C\to E$ has associated graded map equal to the identity
on $C$ and is therefore an isomorphism.
\end{proof}

\begin{corollary}
\label{cor:away-from-two-basis}
Let $A$ be a basic finite-dimensional representation-finite algebra over an
algebraically closed field $k$ with $\charac k\ne2$.
\begin{enumerate}[label={\upshape(\arabic*)}]
\item There is an isomorphism $A\cong k\ray A$, and $A$ is standard.
\item The algebra $A$ admits a normed multiplicative basis consisting of one
representative from each nonzero ray of $A$, with the primitive idempotents
representing the identity rays.
\end{enumerate}
\end{corollary}
\begin{proof}
By Proposition~\ref{prop:graded-realization}(3) and
Lemma~\ref{lem:deletion-brick}, the graded category $\C=\gr_\JJ\ind(\modu A)$
is a finite-dimensional scalar mesh category $k(\Gamma_A)$ satisfying the
brick hypothesis. By Proposition~\ref{prop:mesh-correction-spaces}, its mesh
derivative $D_{q-1}(\C)$ is surjective for every $q\geq3$. Applying
Proposition~\ref{prop:filtered-normal-form} to $E=\ind(\modu A)$ with its
radical filtration $F^n E=\JJ^n$ yields a filtered category isomorphism
$\ind(\modu A)\cong\C\cong k(\Gamma_A)$ fixing the objects, which proves that
$A$ is standard. Restricting this isomorphism to the projective vertices and
taking opposite algebras gives $A\cong A_0$, and combining this with
Proposition~\ref{prop:scalar-normalization} gives
$A\cong A_0\cong k\ray A_0\cong k\ray A$. Finally, transporting the normed
multiplicative ray basis of $A_0$ across this isomorphism yields (2).
\end{proof}

Combining Corollary~\ref{cor:away-from-two-basis} with
Lemma~\ref{lem:nonbasic-basis} completes the proof of Theorem~A outside
characteristic two, as well as Theorem~B. All periodic mesh calculations
required for this branch are given in Appendix~\ref{app:periodic-mesh}; none
of the results of Parts~\ref{part:mesh-deformations}
and~\ref{part:characteristic-two} are needed when $\charac k\ne2$.
